\scanpage{21}

\begin{center}-- 189 --\end{center}
\hfill n$^{\circ}$ 391

En effet, (\S\ 6, n$^{\circ}$ 5, prop. 12), ce groupe quotient est isomorphe au sous-groupe
$G_0$ de $A(R)$ laissant stable un système de racines simples $S(C)$ de $R$ ; il
est clair que tout élément de $G_0$ détermine un automorphisme unique du schéma
de Dynkin de $R$ relatif à $C$ ; mais inversement, un tel automorphisme se prolonge
en un automorphisme unique de $R$ en vertu de la prop. 2, d'où le corollaire.

Les remarques du début et les définitions précédentes montrent que le schéma
de Dynkin d'un système de racines $R$, relatif à un système de racines simples
$(v_i)_{1\leq i\leq\ell}$, détermine le schéma de Coxeter du groupe de Weyl $W(R)$ de la manière
suivante : les graphes sous-jacents sont les mêmes ; en outre, si $\{v_i,v_j\}$
est une arête du schéma de Dynkin, le nombre $m_{ij}$ est égal à 3 si $f(v_i)=f(v_j)$,
à 4 si $f(v_i)/f(v_j)$ est égal à $\frac12$ ou à 2, à 6 si $f(v_i)/f(v_j)$ est égal à $1/3$ ou
à 3. On verra par contre ci-dessous qu'un même schéma de Coxeter peut correspondre
à deux schémas de Dynkin non isomorphes, (en fait, il ne peut correspondre qu'à
\uline{deux} tels schémas au plus, et cela ne se produit que pour les schémas de Coxeter
de type $B_\ell$ ; pour tous les autres types, le schéma de Coxeter de $W(R)$ détermine
le schéma de Dynkin de $R$ à isomorphisme près).

THEOREME 2.-- \uline{Si} $R$ \uline{est un système de racines réduit irréductible}, \uline{son schéma de
Dynkin est nécessairement isomorphe à l'un des suivants}, (\uline{qui sont deux à deux
non isomorphes}) :

$A_\ell$\quad
\diagram{\begin{tikzcd}[column sep=small]
\vert\arrow[r,dash]&\vert\arrow[r,dash]&\vert\arrow[r,dash,dotted]&\vert\arrow[r,dash]&\vert
\end{tikzcd}}
\hfill ($\ell\geq1$ \uline{sommets})

$B_\ell$\quad
\diagram{\begin{tikzcd}[column sep=small]
\overset{2}{\vert}\arrow[r,dash]&\overset{2}{\vert}\arrow[r,dash]&
\overset{2}{\vert}\arrow[r,dash,dotted]&\overset{2}{\vert}\arrow[r,dash]&
\overset{2}{\vert}\arrow[r,dash]&\overset{1}{\vert}
\end{tikzcd}}
\hfill ($\ell\geq2$ \uline{sommets})

$C_\ell$\quad
\diagram{\begin{tikzcd}[column sep=small]
\overset{1}{\vert}\arrow[r,dash]&\overset{1}{\vert}\arrow[r,dash]&
\overset{1}{\vert}\arrow[r,dash,dotted]&\overset{1}{\vert}\arrow[r,dash]&
\overset{1}{\vert}\arrow[r,dash]&\overset{2}{\vert}
\end{tikzcd}}
\hfill ($\ell\geq3$ \uline{sommets})

$D_\ell$\quad
\diagram{\begin{tikzcd}[column sep=small,row sep=tiny]
&&&&&&\backslash\\
\vert\arrow[r,dash]&\vert\arrow[r,dash]&\vert\arrow[r,dash,dotted]&
\vert\arrow[r,dash]&\vert\arrow[r,dash]&
{}\arrow[ur,dash]\arrow[dr,dash]\\
&&&&&&/
\end{tikzcd}}
\hfill ($\ell\geq4$ \uline{sommets})

$E_6$\quad
\diagram{\begin{tikzcd}[column sep=small,row sep=small]
\vert\arrow[r,dash]&\vert\arrow[r,dash]&\vert\arrow[r,dash]\arrow[d,dash]&
\vert\arrow[r,dash]&\vert\\
&&-
\end{tikzcd}}

\mspage{20}
